\begin{proof}
Since $F$ is algebraically closed and nontrivially valued, $|F^\times|$ is
dense in $\mathbf R_{>0}$.  Let
\[
 A' = \{r'\leq|T|\leq R'\},
 \qquad r<r'<R'<R,
 \qquad r',R'\in|F^\times|,
\]
be a strict closed subannulus of $A$.  Both $f|_{A'}$ and its inverse belong
to the annulus algebra of $A'$.  After rescaling the coordinate,
\Cref{prop:ch08-annulus-skeleton} shows that
\[
 s\longmapsto-\log|f(\zeta_{0,\rme^{-s}})|
\]
is concave and piecewise affine with integral slopes on the skeleton of
$A'$.  Applying the same result to $f^{-1}$ shows that its negative is also
concave.  It is therefore affine, with an integral slope.

The maximum inequality in that proposition, first for $f$ and then for
$f^{-1}$, gives
\[
 |f(z)|\leq|f(r_{A'}(z))|,
 \qquad
 |f(z)|^{-1}\leq|f(r_{A'}(z))|^{-1}.
\]
Hence equality holds in both inequalities, so the integral affine formula
on the skeleton gives
\eqref{eq:p062-open-annulus-unit-integral-slope} on all of $A'$.  Such
closed strict subannuli cover $A$, and any two are contained in a third one.
Uniqueness of an affine function on a nondegenerate interval makes their
constants and integral slopes agree.  This proves the lemma.
\end{proof}

\begin{proof}
Fix $\alpha$ and $\beta$.  Every $T_{\beta,k}$ is an invertible analytic
function on $\pi^{-1}(U_{\alpha\beta})$, so
\Cref{lem:ch15-log-smooth} gives
\[
 \mathrm d'\mathrm d''[\log|T_{\beta,k}|]=0.
\]
Use the $\alpha$-trivialization and apply
\Cref{thm:p062-open-polyannulus-products-over-curves} to this function.
Its unique coefficient formula gives a continuous zero-current function
$c_{\alpha\beta,k}$ on $U_{\alpha\beta}$ and locally constant real-valued
functions $M_{\alpha\beta,ki}$ such that
\[
 \log|T_{\beta,k}|
 =\pi^*c_{\alpha\beta,k}
  +\sum_{i=1}^s(M_{\alpha\beta,ki}\circ\pi)
                      \log|T_{\alpha,i}|.
\]
Collecting the $s$ identities gives
\eqref{eq:p062-annular-coordinate-transition}, initially with
$M_{\alpha\beta}$ valued in $\operatorname{Mat}_{s\times s}(\mathbf R)$.

We next prove that every entry is integral.  Fix
$x\in U_{\alpha\beta}$ and a complete algebraically closed extension
 $F/\mathscr H(x)$.  For a fixed $i$, choose $F$-points in all the annular
factors except the $i$-th one; such points exist because $|F^\times|$ is
dense.  Pullback preserves invertibility, so restricting the base change of
$T_{\beta,k}$ to the resulting one-dimensional open annulus gives an
analytic unit $f_{x,i}$.  The coefficient identity is preserved by base
change and restricts to
\[
 \log|f_{x,i}|=\kappa_{x,i}
       +M_{\alpha\beta,ki}(x)\log|T_{\alpha,i}|
\]
for a real constant $\kappa_{x,i}$.  The lemma and uniqueness of the slope
give $M_{\alpha\beta,ki}(x)\in\mathbf Z$.  This holds for every $x$, $i$,
and $k$, so $M_{\alpha\beta}$ is integer-valued.

Apply the same construction with $\alpha$ and $\beta$ reversed and
substitute the two formulas.  The resulting expression of $L_\alpha$ in
the $\alpha$-coordinates has linear coefficient
$M_{\beta\alpha}M_{\alpha\beta}$.  Uniqueness in the cited product theorem,
applied componentwise, compares it with the tautological expression
$L_\alpha=0+I_sL_\alpha$ and gives
\[
 M_{\beta\alpha}M_{\alpha\beta}=I_s,
 \qquad
 c_{\beta\alpha}+M_{\beta\alpha}c_{\alpha\beta}=0.
\]
Both matrices are integer-valued, so every $M_{\alpha\beta}$ lies in
$\operatorname{GL}_s(\mathbf Z)$.  On a triple overlap, substitute the
$\alpha$-to-$\beta$ formula into the $\beta$-to-$\gamma$ formula and compare
with the direct $\alpha$-to-$\gamma$ formula.  The same uniqueness gives
exactly \eqref{eq:p062-annular-affine-cocycle}.

The first identity in \eqref{eq:p062-annular-affine-cocycle} is the
transition law for a local system of lattices.  If the trivializing atlas is
replaced, apply the same construction to the union of the two atlases.  The
integral overlap matrices give an isomorphism between the two resulting
local systems, and uniqueness makes these isomorphisms compatible on triple
overlaps.  This proves the asserted independence up to canonical
isomorphism.  The second identity is the affine cocycle law for the
translation terms after realification.

Now let $h$ have zero current.  Applying the product theorem in each chart
gives the unique expression \eqref{eq:p062-annular-local-normal-form}.
Substitute \eqref{eq:p062-annular-coordinate-transition} into the
$\beta$-expression and compare it with the $\alpha$-expression.  Uniqueness
of the base and slope coefficients gives
\eqref{eq:p062-annular-coefficient-transition}.  The first of those
identities is precisely the dual real transition law, so the $a_\alpha$
define a global locally constant section of
$\mathbb{V}_{\pi,\mathbf R}^\vee$.  The second identity records the remaining
affine translation and completes the proof.
\end{proof}
